% This is samplepaper.tex, a sample chapter demonstrating the
% LLNCS macro package for Springer Computer Science proceedings;
% Version 2.21 of 2022/01/12
%
\documentclass[runningheads]{llncs}
%
\usepackage[T1]{fontenc}
% T1 fonts will be used to generate the final print and online PDFs,
% so please use T1 fonts in your manuscript whenever possible.
% Other font encodings may result in incorrect characters.
%
\usepackage{amsmath}
\usepackage{graphicx}
% Used for displaying a sample figure. If possible, figure files should
% be included in EPS format.
%
% If you use the hyperref package, please uncomment the following two lines
% to display URLs in blue roman font according to Springer's eBook style:
%\usepackage{color}
%\renewcommand\UrlFont{\color{blue}\rmfamily}
%\urlstyle{rm}
%
\begin{document}
%
\title{TM-Vec 2: Accelerated Remote Homology Detection for Structural Similarity}
%
% \titlerunning{Abbreviated paper title}
% If the paper title is too long for the running head, you can set
% an abbreviated paper title here
%
\author{Aryan Keluskar \inst{*,1} \and
Paarth Batra \inst{*,1} \and
Valentyn Bezshapkin\inst{*,3} \and
James T. Morton\inst{2,\textsuperscript{\textdagger}} \and
Qiyun Zhu\inst{1,\textsuperscript{\textdagger}}}
%
\authorrunning{A. Keluskar, P. Batra, V. Bezshapkin, et al.}
% First names are abbreviated in the running head.
% If there are more than two authors, 'et al.' is used.
%
\institute{The Biodesign Institute 
\email{\{akeluska,pbatra6,qiyun.zhu\}@asu.edu}\\ \and
Gutz Analytics
\email{jmorton@gutzanalytics.com}\\ \and
Institute of Microbiology, ETH Zurich 
\email{valentyn.bezshapkin@micro.biol.ethz.ch}}
%
\maketitle            
%
\begin{abstract}
Understanding protein function is an essential aspect of many biological applications. While billions of protein sequences have been identified from sequencing data, the number of protein folds underlying biology is surprisingly limited, likely numbering tens of thousands. The "sequence-fold gap" limits the success of functional annotation methods that rely on sequence homology, especially for newly sequenced, divergent microbial genomes. TM-Vec is a deep learning architecture that can predict TM scores as a metric of structural similarity directly from sequence pairs, bypassing true structural alignment. However, the computational demands of its protein language model (PLM) embeddings create a significant bottleneck for large-scale database searches. To address this, we present two innovations: TM-Vec 2, a new architecture that optimizes the computationally-heavy sequence embedding step, and TM-Vec Student, a highly efficient 2.6M parameter model created by distilling knowledge from TM-Vec 2. Our new models were benchmarked for both accuracy and speed on CATH domains and large-scale database queries. We finally compare it to state-of-the-art models to observe that TM-Vec Student achieves speedups of up to 258x over the original TM-Vec and 56x over Foldseek for large-scale database queries, while maintaining higher accuracy compared to the original TM-Vec.

\keywords{Functional annotation \and
Protein structural similarity \and
TM-score prediction \and
Neural network \and
Knowledge distillation \and
Computational biology}
\end{abstract}

\footnotetext[1]{* denotes equal contribution}
\footnotetext[2]{\textsuperscript{\textdagger} Correspondence: Qiyun Zhu <qiyun.zhu@asu.edu> and James T. Morton <jmorton@gutzanalytics.com> }

\section{Introduction}

Understanding protein function is critical for biological applications, including drug discovery, therapeutic development, and precision medicine. The ability to accurately assess structural relationships between proteins without experimental 3D structures enables applications including remote homology detection \cite{hamamsy2024protein}, functional annotation transfer \cite{koestler2010fact}, and evolutionary analysis \cite{fares2006caps}. Remote homology detection, wherein similarity between proteins exhibiting negligible sequence similarity (\textless 0.25 identity) is identified through conserved three-dimensional architectures, constitutes a fundamental challenge in computational structural biology that has acquired increased relevance following the expansion of predicted structure databases \cite{kilinc2025major}.

The Template Modelling Score (TM-Score) is the gold standard metric for quantifying structural similarity, providing a length-normalised measure that correlates well with fold similarity \cite{zhang2005tm}. 
%\cite{doi:10.1073/pnas.2016239118,elnaggar2021prottranscrackinglanguagelifes}. 
By learning to embed proteins in a space where cosine similarity correlates with TM-scores, TM-Vec \cite{hamamsy2024protein} achieved state-of-the-art performance in remote homology detection. It demonstrated that deep learning approaches can predict structural similarity directly from sequences with remarkable accuracy, something that was previously unattainable using sequence-based alignment techniques. Despite these advances, current deep learning approaches, including TM-Vec, remain computationally intensive, preventing their integration into routine bioinformatic pipelines. The exponential growth of biological sequence databases has generated a critical need for rapid and scalable methods capable of analysing billions of protein sequences efficiently.  

To address this computational bottleneck, we present **TM-Vec 2**, which replaces ProtT5-XL with the efficient Lobster-24M protein language model, reducing embedding costs by 60×. We then introduce **TM-Vec Student**, a distilled model that eliminates PLM dependency entirely by learning directly from raw amino acid sequences using a compact BiLSTM architecture, where cosine similarity between embeddings directly predicts TM-score. **TM-Vec 2 serves as the teacher model** for knowledge distillation, enabling us to train TM-Vec Student—a compact 2.6M parameter model that achieves **258× speedup over the original TM-Vec** while maintaining superior accuracy.

\section{Related Work}


\subsection{Structural Similarity Prediction}

Traditional "gold standard" structural alignment algorithms, such as TM-align \cite{zhang2005tm} and DALI \cite{holm1995dali}, are computationally expensive as they directly superimpose 3D coordinates. This cost was manageable when the PDB was the primary, smaller data source.
However, the release of AlphaFold \cite{jumper2021highly} created a massive data surplus, with hundreds of millions of predicted structures in databases like AlphaFoldDB. This new scale has rendered these methods computationally infeasible for large-scale analysis and created the need for developing novel approaches.
The new tools can be roughly separated into two categories: structure-based and embedding-based. Foldseek \cite{van2024fast} addresses some efficiency concerns through algorithmic optimisations and the 3Di structural alphabet, achieving remarkable speed improvements while maintaining sensitivity comparable to traditional structural aligners. While these methods are highly accurate, they require precomputed or predicted structures, restricting their scalability. 

\subsection{Protein Language Models for Structure}

The emergence of large-scale protein language models has transformed structural biology. ESM \cite{rives2021biological}, ProtBERT \cite{elnaggar2021prottranscrackinglanguagelifes}, and ProtT5 \cite{elnaggar2021prottranscrackinglanguagelifes} demonstrated that self-supervised learning on protein sequences captures evolutionary and, hence, structural information. TM-Vec leveraged this by fine-tuning ProtT5-XL for structural similarity prediction, achieving superior performance compared to sequence-based methods. However, the computational cost of these large models has limited their practical implementation. Recent advances have produced smaller, more efficient protein language models; Lobster \cite{frey2024cramming} achieved competitive performance with orders of magnitude fewer parameters through optimized training procedures. Ankh \cite{elnaggar2023ankh} demonstrated that careful architectural choices can maintain representation quality while reducing computational requirements. Further work has explored various efficiency improvements, such as parameter sharing \cite{sachan2018parameter} and PEFT \cite{sledzieski2024democratizing}, but few studies have addressed the task of efficient structural similarity prediction.

\subsection{Knowledge Distillation in Biology}

Knowledge distillation \cite{hinton2015distilling} has shown promise in compressing large neural networks while preserving performance. In computational biology, DistilProtBERT \cite{geffen2022distilprotbert} successfully reduced ProtBERT's size by 50\% while maintaining accuracy on protein classification tasks. Another paper have showed multi-teacher distillation to reduce time complexity and preserve accuracy across multiple tasks, including binary and mutli-class classification and regression \cite{shang2024accurate}. However, for TM-score prediction, the challenge is not maintaining fine-grained numerical precision, but rather efficiently capturing global structural topology. Since the primary goal is to distinguish between different folds (TM-score > 0.5) rather than minimizing regression error at the decimal level, a massive transformer architecture may be computationally redundant.


\section{Methods}

\subsection{TM-Vec 2: Identifying Efficient Foundation Models}

% \begin{figure}
% \centering
% \includegraphics[width=1\linewidth]{foundation plms.png}
% \caption{\label{fig:foundations} Both of the Lobster models show significant speed-up at inference time compared to our baseline based on ProtT5-XL.}
% \end{figure}

% Improve the diagram by adding notes of lobster's scalability from 2 GPUs to 4 GPUs. Re-draw the diagram in Inkscape for better resolution.

We systematically evaluated three different protein language models as alternatives to ProtT5-XL in the TM-Vec architecture. We examined: 1) Lobster \cite{frey2024cramming} (24M and 150M parameters), an optimized transformer trained with highly efficient attention mechanisms; 2) Ankh    \cite{elnaggar2023ankh}, a larger 450M parameter model featuring architectural optimizations for robust protein sequence encoding; 3) ProtMamba \cite{sgarbossa2025protmamba} (107M parameters), a state-space model that offers a distinct advantage over transformers by achieving linear time complexity with respect to sequence length. 

All inference time benchmarks were conducted using a system with 1x NVIDIA A100 GPU. The experiments utilized a standardized computational environment with identical memory and compute specifications across all model evaluations. Eight distinct sequence length categories were tested: 64, 128, 192, 256, 320, 384, 448, and 512 amino acids. These lengths represent the typical range of protein sequences encountered in structural biology applications. Full results can be found in Appendix. 

\begin{figure}
\centering
\includegraphics[width=0.75\linewidth]{finetuning time.png} \caption{\label{fig:finetuning-time}{Bar graph comparing fine-tuning times for state-of-the-art protein language models.}}
\end{figure}

We also conducted fine-tuning time experiments to identify the foundation model which would be quickest to fine-tune on the TM-Vec training dataset. This was done on a system with 4x A100 GPUs with 10,000 pairs of protein sequences show that Lobster 24M is the fastest to fine-tune, completing in roughly 0.1 hours. Lobster 150M requires slightly more time, hovering around 0.15-0.2 hours, which is still a very fast model to fine-tune compared to our baseline of ProtT5-XL. ProtMamba on the other hand, being a novel architecture with not as much community support, ended up taking more than 6 hours. Therefore, we settled on using Lobster 24M in TM-Vec 2 since it performed well on both inference and fine-tuning speed.

% TODO: Add info from MLSB Appendix A/B here.


 \begin{figure}
 \centering
 \includegraphics[width=\linewidth]{illustration.png} \caption{\label{fig:illustration}{For each protein pair $(X,Y)$, the TM-Vec 2 teacher generates a TM-score. The student model encodes both sequences using a compact BiLSTM architecture into 512-dimensional embeddings, and the cosine similarity between embeddings directly predicts the TM-score. Training minimizes mean squared error between cosine similarity and teacher-predicted TM-scores, enabling the student to learn a metric space where structural similarity is captured by embedding geometry.}}
 \end{figure}


\subsection{TM-Vec Student: Knowledge Distillation Architecture}

The TM-Vec Student model eliminates PLM dependency by distilling knowledge from TM-Vec 2. Operating directly on raw amino acid sequences, the student comprises 2.6 million parameters—a 1,000× reduction compared to the ProtT5-XL architecture. A key insight driving our design is that TM-score prediction can be reformulated as a metric learning problem: rather than learning a regression function on fused embeddings, we train the encoder to produce representations where cosine similarity directly approximates the TM-score.

\subsubsection{Sequence Encoder}
The architecture utilizes a BiLSTM backbone \cite{hochreiter1997long,graves2005framewise} prioritizing linear space complexity over the quadratic costs of transformer-based alternatives \cite{abduljabbar2021development}. The encoder maps each amino acid to a 128-dimensional embedding, followed by a 2-layer BiLSTM with 512 hidden units (256 per direction). An attention pooling mechanism aggregates the bidirectional hidden states into a single 512-dimensional sequence representation:
\begin{equation}
\alpha_t = \frac{\exp(w^\top h_t)}{\sum_{t'} \exp(w^\top h_{t'})}, \quad z = \sum_t \alpha_t h_t
\end{equation}
where $h_t$ is the concatenated bidirectional hidden state at position $t$, and $w$ is a learned attention vector. The pooled representation passes through a linear projection followed by layer normalization to produce the final embedding $z \in \mathbb{R}^{512}$.

\subsubsection{Cosine Similarity as TM-Score Prediction}
Unlike traditional knowledge distillation approaches that employ feature fusion and regression heads, TM-Vec Student directly uses the cosine similarity between sequence embeddings as the predicted TM-score:
\begin{equation}
\widehat{\text{TM}}(A, B) = \cos(z_A, z_B) = \frac{z_A \cdot z_B}{\|z_A\| \|z_B\|}
\end{equation}
This design eliminates the need for a dedicated predictor network, reducing both parameter count and inference latency. The model learns to position proteins in the embedding space such that structurally similar proteins (high TM-score) have high cosine similarity, while dissimilar proteins are placed apart.

\subsubsection{Training Objective}
The student model is trained with a simple mean squared error loss that directly aligns cosine similarity with TM-scores predicted by the TM-Vec 2 teacher:
\begin{equation}
\mathcal{L} = \frac{1}{N} \sum_{i=1}^{N} \left( \cos(z_{A_i}, z_{B_i}) - \text{TM}_i \right)^2
\end{equation}
where $\text{TM}_i$ is the **TM-Vec 2 teacher-predicted** TM-score for protein pair $(A_i, B_i)$. Distilling from TM-Vec 2 (rather than the original TM-Vec) improved student validation $R^2$ from 0.96 to 0.983, justifying our two-stage approach. This formulation implicitly handles the bounded nature of both cosine similarity $[-1, 1]$ and TM-scores $[0, 1]$---the model naturally learns to produce non-negative cosine similarities for protein pairs, as TM-scores are always positive. The simplicity of this objective, compared to multi-component loss functions, proved sufficient for achieving high accuracy while enabling stable training dynamics. 

\section{Experimental Setup}

\subsection{Optimized TM-Vec 2 Training} %TODO: Change name

We conducted training for the new TM-Vec 2 model on NAIRR resources provided by the NCSA's DeltaAI cluster. To ensure high representation of structurally similar proteins,the 146M pair dataset was adjusted to include a 5/90/5 split across three TM-score bins: 0-0.17 (low similarity), 0.17-0.6 (medium similarity), and 0.6+ (high similarity). The model trained on 1x GH200 node for 36 hours, and saw 10.6M samples per epoch, balanced according to the bin-based splitting strategy.

\subsection{TM-Vec Student Training Configuration}

The TM-Vec Student model was trained on three million protein pairs with TM-scores generated by the TM-Vec 2 teacher model. The dataset was split into 2.55M training pairs and 450K validation pairs (85/15 split). Training was conducted for 25 epochs on a system with 1x NVIDIA A100 GPU for approximately 30 hours, using a batch size of 64, AdamW optimizer with learning rate 1e-3, weight decay 0.01, and OneCycleLR scheduling. The model achieved a final validation $R^2$ of 0.983 and MAE of 0.015, demonstrating that the simple cosine-similarity formulation effectively captures the structural similarity relationships learned by the teacher.

\subsection{Benchmark Datasets}
We evaluated all of the models on the CATH S100 dataset, a non-redundant subset of the CATH protein domain database \cite{pearl2003cath}. Specifically, we collected the first 1,000 protein domains from the CATH S100 dataset. They represent four classes, 26 architectures, 223 topologies, and 328 homologous superfamilies. Benchmarking was performed using a computer system with an AMD EPYC 7413 CPU and an NVIDIA A100 GPU (80GB VRAM).

\section{Results}

\subsection{Efficiency Gains}

\begin{table}[htbp]
\centering
\resizebox{\textwidth}{!}{%
\begin{tabular}{|cc|ccccc|}
\hline
\textbf{Query Size} & \textbf{Database Size} & \textbf{Blast} & \textbf{TM-Vec} & \textbf{Diamond} & \textbf{Foldseek} & \textbf{TM-Vec Student} \\
\hline
10   & 1000   & 9    & 0.6  & 0.521 & 0.640 & \textbf{0.087} \\
100  & 1000   & 16   & 5    & 0.512 & 1.785 & \textbf{0.096} \\
1000 & 1000   & 65   & 50   & 0.597 & 10.934 & \textbf{0.194} \\
10   & 10000  & 10   & 0.7  & 0.573 & 0.612 & \textbf{0.127} \\
100  & 10000  & 25   & 5.2  & 0.644 & 2.648 & \textbf{0.114} \\
1000 & 10000  & 131  & 50.2 & 0.775 & 18.640 & \textbf{0.336} \\
10   & 100000 & 14   & 0.7  & 1.234 & 0.721 & \textbf{0.162} \\
100  & 100000 & 63   & 5.3  & 1.395 & 2.981 & \textbf{0.272} \\
1000 & 100000 & 475  & 51   & 1.743 & 20.832 & \textbf{1.692} \\
\hline
\end{tabular}%
}
\vspace{2pt}
\caption{Search and alignment speed for TM-Vec Student Model is compared against BLAST \cite{altschul1990basic}, TM-Vec, Foldseek and DIAMOND \cite{buchfink2021sensitive} for different query and database sizes.}
\label{tab:query_benchmark}
\end{table}



Encoding time benchmark (Table \ref{tab:encoding_benchmark}) and Query time benchmark (Table \ref{tab:query_benchmark}), as obtained from TM-Vec \cite{hamamsy2024protein}, show substantial improvements in efficiency over TM-Vec and Foldseek. We use a precomputed database of protein structures obtained from the RCSB Protein Data Bank \cite{rose2016rcsb} to benchmark Foldseek, while the TM-Vec models directly encode and query with protein sequences. The TM-Vec Student model requires $0.967$ seconds for encoding 10000 sequences compared to TM-Vec's $625.72$ seconds and Foldseek's $91.446$ seconds. The student model, by eliminating the dependency on pLMs and using simple feed-forward operations, outperforms Foldseek and shows a 60 to 600× improvement in encoding efficiency over TM-Vec while maintaining comparable accuracy. We observe similar speed-ups while querying against databases of various sizes, as documented in Table \ref{tab:query_benchmark}.



\subsection{Prediction Accuracy}

\begin{figure}
    \centering
    \includegraphics[width=1\linewidth]{hexbin_results.png}
    \caption{Pearson R Comparison of Foldseek, TM-Vec 1, TM-Vec 2, and TM-Vec Student} (In Order). TM-Vec models predict pairs more confidently at low similarity as compared to Foldseek, yet maintain similar accuracy.
    {\label{fig:hexbin-results}}
\end{figure}

When comparing TM-Vec results to SoTA, accuracy values are consistent with results generated from Foldseek and correlate well with TM-align results.  

\begin{credits}
\subsubsection{\ackname} This research was supported by the U.S. Department of Energy, Office of Science under award DE-SC0024320. This work utilized computational resources of the DeltaAI supercomputer at the University of Illinois Urbana-Champaign, through the National Artificial Intelligence Research Resource (NAIRR) Pilot program (award: NAIRR250131), and computational resources of the Sol supercomputer at Arizona State University.

\subsubsection{\discintname} The authors have no competing interests to declare that are relevant to this work.
\end{credits}

\bibliographystyle{splncs04}
\bibliography{ref}

\appendix

% \clearpage

% \section{Student Model Architecture}

\begin{figure}
    \centering
    \includegraphics[width=0.65\linewidth]{tmvec student.png}
    \caption{Student Model Architecture}
    \label{fig:architecture}
\end{figure}

\end{document}